\section{Simulation Workflow}
\label{sec:simulation_workflow}

The previous sections introduced the architectural layers of \syssimx{} individually.
This section shows how they come together in a typical user workflow.
Listing~\ref{lst:workflow} illustrates the complete path from component instantiation to result retrieval.

\begin{lstlisting}[
  caption={Minimal workflow for defining and running a co-simulation 
  in \syssimx{}.},
  label={lst:workflow},
  language=Python]
from syssimx import System, Connection
from syssimx.algorithms import GaussSeidelAlgorithm

# 1. Instantiate components (each implements CoSimComponent)
source = LinearSource(name="Source", a=1.0, b=0.0)
integrator = Integrator(name="Integrator", x0=0.0)

# 2. Assemble system and connect interfaces
system = System(name="QuickstartSystem")
system.add_component(source)
system.add_component(integrator)
system.add_connection(Connection(
    src_comp="Source", src_port="y",
    dst_comp="Integrator", dst_port="u",
))

# 3. Select algorithm, initialize, and run
system.set_algorithm(GaussSeidelAlgorithm())
system.initialize(t0=0.0)
system.run(t0=0.0, tf=5.0, dt=0.1)

# 4. Retrieve simulation results
history = system.get_history()
\end{lstlisting}

The listing follows four phases: component instantiation, system assembly 
with explicit connections, execution with a user-selected algorithm, and 
result retrieval.
Both \texttt{LinearSource} and \texttt{Integrator} implement the 
\texttt{CoSimComponent} interface introduced in 
Section~\ref{sec:component_interface}; any component that fulfills this 
contract can participate in the workflow without modification.
The \texttt{System} handles dependency analysis, execution-order 
computation, and connection validation transparently during initialization.
